\section*{4. \textit{Gnomes \& Goblins}}

\subsection*{The Tutorial: Coming to Dwell}

\subsubsection*{The Candle: A Threshold Crossed by Discovery}

The player begins in the opening menu — an open, grassy expanse that extends into the horizon — confronted by several perception-signs: a signpost, each of its directions hung with an unlit candle, and a single lit candle resting on a stump nearby. This is the only part of the program that offers any instruction at all; and it is easily missed, for it shows only if the player thinks to look at the tracked hands, where a line of text says to take up the lit candle and set its flame to one of the unlit ones that adorn the signpost. Even that much resists: the lighting is ``more difficult than one might imagine; it takes some experimenting'' with the candle before the flame catches and the way opens (Salvo, \textit{Veri(dis)similitude}, 31).

What the fumbling teaches is a meaning as much as a grip. At first the signpost reads as a signpost — a trail-sign, a ``direction tone,'' pointing the ways one might go — and the candles as things that merely burn; only in the handling do they take on another sense, an ``interface tone,'' a thing the player must work in order to begin the experience at all (Salvo, \textit{Veri(dis)similitude}, 31). And nothing announces the turn. The one instruction names the task and not its meaning, so that the shift from sign to interface is left for the player to author, out of the player's own trying. Here, in little, is the method the whole world will run on: a sense withheld, and then drawn from the player's own image-making (\textit{phantasia}, \gk{φαντασία}) rather than declared.

Nothing in this is yet a dwelling. The player stands at the edge of the world and not within it — the first skill of grip and reach only beginning to take hold, no surround yet closed about the player, and nothing yet mattering. What the candle lays down is just the disposition the rest will build upon and the rupture will, in the end, turn: a habit of finding the world's sense by working at it, cued by almost nothing. The flame catches, the candle on the direction marked ``Goblins'' takes light, and the menu gives way to a forest.

\subsubsection*{The Forest: Dwelling, and What Withholds}

Lighting the candle on the direction marked ``Goblins'' transports the player to the main world: a fairy-like forest, complete with tiny homes in hollowed-out trees and rope suspension bridges strung between the dwellings, fireflies buzzing in the dark. And the player is not before this world but within it — set down where a moment ago there was a room. This is the designed surround the reading will keep returning to, a virtual \textit{Umwelt} in Uexküll's sense: a surrounding world closed about the one who lives in it like a soap-bubble, furnished to be inhabited and not surveyed (Salvo, \textit{Veri(dis)similitude}, 27). And in virtual reality the inhabiting is nearly literal. Physical locomotion in the program is one to one — a single step forward is just that, a single step — so a real stride becomes a stride among the trees and a real reach a reach toward them, the tracked hands the player's own, no avatar interposed (Salvo, \textit{Veri(dis)similitude}, 28). The spatial and bodily floor of any dwelling — the felt sense of being somewhere, able to move and to handle — is here not conscripted from the player's imagination but supplied by the medium. The cornerstone of the dwelling is set in the body.

For a time the world simply obliges. The player is tempted to investigate the tiny treehouses, invited to push open their wooden shutters and peek inside; sticks and bridges sway when pushed, a fallen fruit can be lifted, an acorn taken up. The forest answers the player's dealings smoothly and at once, and a fluency builds — each handled thing confirming the last, closing again and again the circuit of perceiving-and-doing that Uexküll called the functional circle, until the movements ask no thought. This is a kind of kinesthetic learning, and it is how the dwelling is laid down: not declared but accreted, the player interiorizing the world's mechanics until they are second nature, a skill thickening into the settled ease of being at home somewhere (Salvo, \textit{Veri(dis)similitude}, 30).

One thing, though, refuses. A curiously shaped rock lies on the ground, and the hand that reaches to take it closes on nothing it can use: the rock is, in the earlier study's words, ``resistant or entirely unserviceable'' (Salvo, \textit{Veri(dis)similitude}, 40). It cannot be lifted or moved or turned to any task; and in refusing, it shows what the obliging world had kept hidden. The compliant objects had been ready-to-hand, so transparent in use that they drew no notice; the rock, balking, goes un-ready-to-hand in Heidegger's sense — the thing that announces itself precisely by breaking down under the hand that reaches for it (Salvo, \textit{Veri(dis)similitude}, 40). What balks is then set aside: the rock falls out of the player's dealings and becomes a mere perception-image, a thing that, proving unusable, ceases to matter.

What refuses in this way is, in Rickert's sense, a thing and not an object — something whose being is not used up in what can be done with it, exceeding the use-relation and withdrawing into a depth no dealing reaches (Rickert, \textit{Ambient Rhetoric}, 199–201). Almost everything else in the forest hides that withdrawal by complying; the rock, in refusing, makes it briefly palpable. And yet open refusal is not this world's main way with the player. The forest dwells the player less by pushing back than by holding back — withholding cues, withholding instruction, withholding the meaning of its candles and the consequence, soon, of its bell. The old adage ``less is more'' is at work in the design: by limiting the interactive objects and saying almost nothing, it turns the player's dealings onto the world itself and leaves the player to supply what is kept from view (Salvo, \textit{Veri(dis)similitude}, 33). The rock withdraws from use; the design withholds the world's sense; and both throw the player onto the same conscripted image-making. What looks like a generous world is in this one respect austere: it hands over its objects readily and its meanings almost never.

Into this dwelling a first relation is seeded. A goblin appears — shy, glimpsed at the edges, then nearer. So the dwelling stands: a body at home in a tracked space, a fluency that needs no thought, a world of things that mostly answer and one that would not, a creature met.

\subsubsection*{The Fruit: How a World Moves the Player Without a Word}

Of all the moments in the opening, the fruit exchange is the one in which the design's method is most nakedly itself. It moves the player to do exactly what it wants — pick up a fruit and give it to a goblin — without a single instruction, and it does so by arranging a world rather than issuing a command. Walked through the actualization of desire step by step, the curation becomes legible; and because nothing here is carried by narration or told outright, the scene is the cleanest demonstration the artifact offers of environmental rhetoric at work.

The chain begins before the fruit falls, at the horizon (A₀). By this point the player has had some minutes of unguided play — opening shutters, lifting acorns and sticks, pushing branches — and out of that play a disposition has quietly settled: in this world, things are to be picked up. The forest has also, over these same minutes, produced a goblin; the shy creature makes only brief appearances at first, always on the periphery, glimpsed and gone (Salvo, \textit{Veri(dis)similitude}, 32), but it has drawn the player's attention, so that the horizon now holds not only a field of graspable things but a creature the player has begun to attend to. Neither disposition was taught; both were cultivated by the design, and both are about to be drawn upon.

The trigger is a sound. High in a tree the goblin turns to a piece of fruit on one of the higher branches; grunting with effort and unable to reach it, it goes briefly inside to sharpen its axe, then returns to cut the fruit from the branch (Salvo, \textit{Veri(dis)similitude}, 32). The player hears the chopping before seeing it — a sound the headset anchors in space, turning the head toward its source. This is the first and subtlest stroke of curation: attention itself is directed, at the level of perception (A₁), by an authored cue, so that the player looks up in time to watch the fruit come loose and fall to the moss (06:52.8). What the perception lays down is not yet a plan but a charge — a thing has just come free, and it has come to rest at the player's own level.

From that charged perception the player forms the image of a possible act (A₂): \textit{a fruit I can pick up}. The image is under-determined — nothing has said what the fruit is for — but it is not conjured from nothing. It is filled out from the disposition the free play laid down, the settled sense that the graspable thing is to be grasped. The design has withheld all instruction; and into that withholding the player's own image-making (\textit{phantasia}, \gk{φαντασία}) supplies the act the design declines to name.

What turns \textit{pick it up} into \textit{give it to the goblin} is the goblin itself. Once the fruit is down, the creature waits: it merely looks at the fruit and then at the player, no word exchanged, and it is left to the player to figure out what it wants (Salvo, \textit{Veri(dis)similitude}, 32). That bare expectant nearness is the second authored cue, and on it the player ratifies a settled judgment (\textit{doxa}, \gk{δόξα}): this just fell, and the creature wants it (A₃). No line of text proposes the gift; the goblin's posture, read through the player's own \textit{phantasia}, does. The mood that colours the judgment is the helping tone the encounter has set — for the goblin takes its meaning from the player's own comportment, an abrupt approach reading as a threat and sending it scampering, a gentle one drawing it near, so that curiosity tips over into care (Salvo, \textit{Veri(dis)similitude}, 35–36).

And so the player acts (A₄): carries the fruit to the goblin and sets it on the ground between them (07:13.3). The act is the player's own — a voluntary motion sourced in a curiosity become helpfulness, a desiderative (\textit{epithymetic}) impulse completing itself as a wish (a bouletic \textit{orexis}) to help, its means, the offering, deliberately chosen. Nothing forces it; it is no forced motion (\textit{bia}, \gk{βία}). The freedom is real. But the freedom lies in the engaging, not in the outcome: the design has so arranged the scene — the chop, the fall to eye level, the patient wait — that of everything the player might do, this is the act made salient, made likely, very nearly made inevitable. Here is the whole principle of the thing — curation, not coercion: the player does just what the designer intends, and does it freely.

The goblin's part seals the reading. It does not take the fruit at once; it waits some eight seconds, then bends, lifts it, and carries it off (07:21.7). That delay is the tell. The goblin is not a fellow agent answering a gift but a scripted trigger firing on a timer — programmed responsiveness, the efficient causation of the hidden agent, the programmer, for which the goblin's ``wanting'' stands in. In the earlier study's plainer words, the creature's actions ``are nothing more than… programmed responses triggered when the user does something the program recognizes,'' so that even as it appears to act on its own, it does so only as a superficial agent whose motive is not its own but the programmer's (Salvo, \textit{Veri(dis)similitude}, 36). The player reads none of this: \textit{phantasia} fills the pause as expectancy and the script as intention, so that a timed retrieval is lived as a creature's gratitude.

What the exchange leaves behind (A₄→A₀′) is small and real: a helping tie where there was none, and with it the first sediment of an ethical disposition (\textit{hexis}, \gk{ἕξις}) toward the goblin — the player is now one who has helped, and the world one in which help is taken up and answered. The Shared dimension of involvement, dormant until now, has opened; the dwelling has gained a second person. It is this seeded relation that the later beats presuppose and deepen — the descent among the goblins, and, before it, the rupture the whole dwelling makes possible.

Read whole, the fruit exchange is the persuasive virtual ecology in miniature. Every cause but two is the designer's: the fruit, the goblin, the sound, the fall, the wait, and the rule that binds them. The two that remain the player's — the desire that moves the act and the end it serves — are exactly the ones the design cannot author and has no need to, because it has built a world in which the player's desire and the designer's end fall into a single line. To curate an environment, on this evidence, is to arrange the conditions under which a free act becomes the likely one: to move someone to do what you would have them do by handing them a world, and never a word.

\bigskip

\subsection*{The Disruption: Rupture and Carry-Over}

\subsubsection*{The Miniaturization: A Rupture That Re-Incorporates}

By the time the gold bell lies on the moss at the tree-base, the player has already come to dwell. The opening minutes do their quiet work: a fluent grip on the controls, a path learned by walking it, a goblin met and handed a fallen fruit — a skill (\textit{technē}) settled into a disposition (\textit{hexis}, \gk{ἕξις}) and a relationship begun, sedimented act by act, until the forest is no longer a scene surveyed but a place inhabited. It is into this dwelling that the bell is rung; and what the ringing does turns on its having been rung from within a world already the player's own.

The act itself is small, and entirely the player's. Curiosity moves it — what does this bell do? — a desire (\textit{epithymia}, \gk{ἐπιθυμία}) sourced in the player and in no goal the design has set; and the means is chosen, the bell raised and sounded as a deliberate probe. The ringing is thus no \textit{bouletic} act sourced in some end the world has set, but a \textit{prohairetic} one: a bare \textit{epithymetic} impulse — curiosity — completing itself by a deliberately chosen means, which is what \textit{prohairesis} is, the deliberative completion of a wish and not a fourth kind of motive beside the three desires. What the player cannot author is the consequence. The design withholds it: nothing on the bell, nothing in the wordless world, declares what ringing will bring. So the charged image the player acts from — the \textit{phantasma} (\gk{φάντασμα}) of a bell to be rung — is under-determined as to its outcome, filled out from the player's own settled judgment (\textit{doxa}, \gk{δόξα}) that bells make a pleasant noise. The freedom is real but placed: it lies in the \textit{whether}, the choosing to ring, not in the \textit{what-it-does}, which is foreclosed.

Then the world overwrites the image. Fireflies wake and swirl, light saturates to white, and the point of view drops as the world swells — grass risen into trees, the human scale of a moment before collapsed to the scale of an insect (07:48–08:14). For this the player has no fore-structure; the \textit{phantasma} rung from is exceeded, not fulfilled. And yet the rupture does not expel. The player is transformed, not dragged from the agent's place to the patient's — movement, gaze, a second ringing all remain — and the change runs with the player's desire rather than against it, and so it is no forced motion (\textit{bia}, \gk{βία}), no cutscene seizing the body, but a designed and productive rupture. What it ruptures is the scale at which the player was incorporated; what it leaves intact is the incorporation itself.

It is the deepest of several breakdowns the world can spring, and the only one that re-incorporates. A thing had already withdrawn from use — the rock that would not be handled — and disclosed, in withdrawing, that the world need not yield; the apparatus itself can fail, the headset slipping or the power cutting, and disclose the rendered world as rendered by tearing it away — in that failure one event befalls the digital object the machine still holds and another befalls the virtual world the player inhabited, which simply ceases; even a held image can flip, gestalt-like, and unsettle perception for a beat. Each is a rupture that discloses by breaking. The miniaturization breaks more deeply than any of them — it remakes the very scale of being-in-the-world — and yet, alone among them, it breaks without expelling: the disclosure it works is undergone from within a dwelling it does not end.

This is why the shock lands as wonder and not as dread. The mood the forest had set — a childish, unguarded delight — colours the perceptual moment before the transformation arrives, so that when the world gives way the emotion concretized is anxiety in its disclosive register and not its threatening one: the vertigo of a ground withdrawn, recovered almost at once into the poise of an agent still standing. Anxiety (\textit{Angst}) here individualizes without unhoming for good.

Rendered as motive, the episode runs on a single ratio. The player attributes the change to the bell, to the world — \textit{the scene did it} — and the scene, in transforming, transforms what the agent is and can do: a Scene–Act relation, the container reshaping the contained (Burke, \textit{Grammar}, 443). Two terms go concealed in that attribution, as a banished term goes on doing its work unseen (Burke, \textit{Grammar}, 21): the player's own settled readiness to ring intriguing things, and the maker who wrote the rule — the programmer, the hidden Agent for whose efficient causation the magic stands in.

Read as union, the beat is the height of the whole arc. Everything the forest had built — the fluent body, the inhabited space, the wonder, the dawning story — is co-active at once, and incorporation reaches the apex of its climb from off-line anticipation to absorbed play (Calleja, \textit{In-Game}, 178). The rupture is then the limit case of the strictest condition: an intrusion that does not break the absorption but is taken up into it (Calleja, \textit{In-Game}, 172), the porosity by which the world's most violent move deepens rather than dispels belonging — the dissimilar shock re-incorporating the player at a new scale instead of casting the player out. World-disclosedness (\textit{Welterschlossenheit}) — for Heidegger, disclosed equiprimordially with the others who share the world and with one's own existence (Heidegger, \textit{Being and Time}, 176) — is not undone but re-pitched: the designed surround is disclosed again, smaller and stranger, and again as a place to be dwelt in.

This is \textit{veri(dis)similitude} made event. A world verisimilar enough to be entered through fluent dwelling is broken open by a possibility dissimilar enough that the breaking carries ontological force; and the force is disclosive. What the miniaturization lays bare — what it \textit{unconceals} — is the contingency of a constitution that ordinarily withdraws from notice: that the player's scale, orientation, and standing in a world are not given but rendered, and can be rendered otherwise. The bell that worked the rupture becomes, now, a way through: rung again, it teleports, and the player moves by it among the small world's places — the threshold at which world-disclosedness turns toward the others who dwell there.

\subsubsection*{The Tavern: Co-Disclosedness at Its Threshold}

Rung from within the miniaturized world, the bell no longer transforms; it carries. Each ringing sets the player down in a new place among the tiny world's locations, and one of these is a tavern — a drinking-hall built for a crowd, its long tables and benches and scattered mugs sized for company, lamplit and warm. Two goblins are at home in it: one asleep, one bent over a drink (08:56–09:16). Nothing here is staged for the player's coming. The room and its dwellers go on with a life of their own; and what the place discloses, in disclosing that life, is the goblins' being-with-one-another — a way of living together already under way before the player arrived among it, and not pausing for the arrival.

To be set down in that hall is to undergo one of the founding disclosures of being-in-a-world. The world the player now inhabits is a with-world: its dwellers come forward not as objects met but as others already living together, so that to enter their place is to receive the ``disclosedness of the Dasein-with of Others'' that belongs to every being-with (Heidegger, \textit{Being and Time}, 160). And the disclosure runs both ways. At some of the bell's destinations a goblin lifts its head and looks — a small thing, and the whole of it: in the returned look the other is given as an other, and the player as one who is there to be seen, co-present in a shared place. Co-disclosedness (\textit{Miterschlossenheit}) opens here not as a doctrine but as this mutual look across a lamplit room.

Read as motive, the tavern works by what Burke names consubstantiation — a shared substance built out of shared act. ``A way of life is an acting-together,'' and in acting together the members come to hold the sensations and attitudes that make them of one stuff (Burke, \textit{Rhetoric}, 21). The goblins' communal drinking is just such an acting-together, their consubstantial way of life made visible; and the warmth of the inhabited scene works on the player as the sympathetic imagination works, rendering the goblin-world present as something that matters, something to belong to, before any deed is done among them (Burke, \textit{Rhetoric}, 83–84). What the player cannot do is enter the substance. Admitted to the hall and co-present in it, the player remains ``both joined and separate'' (Burke, \textit{Rhetoric}, 21) in the thinnest register the phrase will bear: joined as a guest is joined, present and unturned-away; separate as a stranger is separate, unaddressed, no member of the we that drinks.

For the joining stops at the look. The goblins drink and sleep and raise their heads, but they never act with the player — never hand a cup, never make room at the table, never take up a common task. Their motion is programmed responsiveness, scripted by the hidden maker whose efficient causation the world's life stands in for; and what fills it out as a life — the sense that someone is genuinely at home behind the drinking — is supplied by the player's own image-making (\textit{phantasia}, \gk{φαντασία}). The goblins hold no settled disposition (\textit{hexis}, \gk{ἕξις}) of their own to extend, so nothing in them can close upon the player as a fellow. The acting-together is given to be witnessed and not to be entered. What the design discloses, then, is the co-world entire — the shared dwelling, the others at home in it, the look that grants co-presence — and what it withholds is the single thing that would consummate the bond: a deed done together. The consubstantial ``we'' is not merely nascent here, waiting on more time; it is structurally held back, designed out — whether by deliberate authorial choice or by the limits of the NPC system, the analytic result is one and the same — shown as a thing the player may stand inside of and yet have no part in. This is co-disclosedness given as a preview — the disclosure of being-with isolated, deliberately, from its enactment.

So the small world holds the player at two depths at once. World-disclosedness is sustained and even widened: the bell that re-scaled the surround now ferries the player through it at will, and the dwelling deepens with each place entered. Co-disclosedness arrives, but only to its threshold — the co-world and the mutual look disclosed, the acting-together kept just beyond reach. These are not two unions but one incorporation disclosing along its two poles — one in substrate, different in being; and \textit{Gnomes \& Goblins}, by opening the second pole exactly to the edge of consummation and no further, lets that single incorporation be seen as a consummated collective never could. That the co-world shows at all, even with its enactment withheld, is the work of the same disclosive power the rupture exercised in the breaking — the power of the rendered world to make briefly visible what it silently constitutes.

\subsubsection*{Carry-Over: The Unconcealing and What It Leaves Behind}

The rupture's deepest work is not the change of scale but the change of sight. A world dwelt in withdraws from notice; the fluent body, the inhabited place, the standing from which one acts do their work precisely by going unremarked — the hammer disappearing into the hammering. What the miniaturization does is break that transparency for an instant, and in breaking it, show it. The player who has come to dwell at human scale is handed, without warning, a world in which grass stands like trees; and the gift is not the new scale but the sudden visibility of the old one \textit{as} a scale — one rendering among possible others, contingent where it had seemed simply given. This is an unconcealing: the breakdown lays bare the ordinarily-hidden constitution of being-in-a-world, the way a tool laid bare in its failing discloses the whole equipmental web that had silently carried it.

The mood is the proof. Anxiety (\textit{Angst}) is not here a failure of absorption but its disclosive limit — the affect in which the player's enworldment, having lost its footing, comes into view \textit{as} enworldment. Anxiety individualizes: it draws the one who undergoes it out of the absorbed, anonymous everyday and sets the bare fact of being-in-a-world before the player. The miniaturization works that individualizing stroke in little, and works it as wonder rather than dread; but the structure holds. What had been lived is, for a moment, seen.

This recasts the union itself. To be incorporated into a designed world is, by default, to be absorbed in it uncritically — a dwelling at once genuine and fallen. The fallenness is precise: the world's significance comes pre-articulated, settled in advance as the anonymous ``they'' settles the sense of the everyday, so that the player takes on trust a meaning another has authored; the player is drawn ever onward by curiosity, which for Heidegger is less a thirst for the new than a fleeing of oneself into whatever distracts; and beneath both, the player is constituted as a functional object of the programmer's design, moved through a possibility-space foreclosed in advance — every choice, like the chord already lying in the piano before the child's hand finds it, available only because it was built in (Adorno, \textit{Aesthetic Theory}, 32). Incorporation left to itself is this fallen union. The rupture is what can convert it: by making the constitution visible, the unconcealing opens the possibility of a dwelling no longer merely undergone but owned — a critical, individualized inhabiting in place of an absorbed one. The productive rupture is thus not the enemy of incorporation but its completion: the point at which a union that was only suffered becomes one that can be answered for.

And the disclosure does not end when the headset comes off. A single first-hand datum marks the carry-over: after several hours in the rendered world, the return to the unrendered one is not immediate — for a quarter-hour or so the actual world wears the strangeness the virtual one had shed, the eye reaching for a scale no longer there, the body carrying an orientation the room will not answer. The report behind that description is first-hand and exact: the return brought ``the odd sensation that I remained tethered to the computer,'' depth-perception briefly unreliable, and the disturbance held until ``it took around 15 minutes for my phenomenological and ontological orientation to recalibrate to my `Being-in-the-world' of objective reality'' (Salvo, \textit{Veri(dis)similitude}, 58 n. 15). The disposition sedimented across hours of dwelling carries over the threshold of the exit and presses, for a time, upon the disposition it had left behind. This is the actualization chain closing its longest loop: the completed act settling into a standing disposition (\textit{hexis}, \gk{ἕξις}), and the disposition outlasting the world that laid it down.

The verdict the evidence will bear is partial but real — a transient supplanting of the player's dispositional ground, an alteration of settled orientation deep enough to bend how the actual world is met on return. It must be held within a limit. The virtual is distinct in kind from the non-virtual; no depth of carry-over converts the rendered world into an unrendered one, and no recalibration makes the player a new kind of being. What the carry-over alters is the \textit{how} of a dwelling, not the \textit{what} of a substance — an alteration (\textit{alloiōsis}, \gk{ἀλλοίωσις}) so deep it strains toward a remaking of character and yet, the virtual keeping its kind, cannot become one. The supplanting inflects the disposition; it does not author a new soul.

Here the method earns its stake. If a designed world can drive the actualization of desire to a substantial union, and if that union sediments into a disposition that carries past the world's edge, then to design such worlds is to design dispositions — a rhetoric that works not on belief but on the standing from which belief and action set out. The miniaturization is a small and benign instance: a moment of wonder that, by unconcealing the contingency of a place, leaves the one who undergoes it faintly re-pitched toward the world. The instance is small; the principle is not. It is the warrant for taking the rendered world seriously as a maker of who its inhabitants are becoming — and the charge under which the design of such worlds must be thought.

\bigskip

